% ============================================================================
% Fichier  : partie0/P00_scene_numerique.tex : FORN (LPI342)
% ============================================================================

\chapter{Prologue : la Scène Numérique}
\label{chap:prologue}
\chaptermark{Prologue}

\epigraph{%
    « Tout contact laisse une trace. »%
}{Edmond Locard, 1920}

\niveauChapitre{Introductif}{Aucun prérequis. Les notions sont reprises en
détail dans les chapitres.}

\begin{boxobjectifs}
\begin{itemize}
    \item Découvrir la posture de l'enquêteur numérique.
    \item Distinguer l'enquête de l'attaque, la trace du jugement.
    \item Découvrir le fil rouge LiZbiZ et l'opération TSANAGA.
    \item Situer le Trilemme CRO et la chaîne de possession.
    \item Formuler les questions auxquelles le cours répond.
\end{itemize}
\end{boxobjectifs}

\bigskip

% ============================================================================
\section*{Un ordinateur et une clé USB, scellés dans un sac}
% ============================================================================

Après le rançongiciel de mars, LiZbiZ découvre qu'un ancien employé de son
site de Garoua, M.~NDONGO, aurait, des mois plus tôt, copié puis revendu le
fichier des clients. Le 20~avril 2026, sur réquisition du parquet, un
ordinateur portable et une clé USB sont saisis à son domicile et remis, sous
scellé, au cabinet \sigle{NYONG SÉCURITÉ}, désigné comme expert.

\medskip
Sur ces supports se trouvent, peut-être, la preuve de ce qui s'est passé.
Mais une preuve mal recueillie, mal conservée, ne vaut rien devant un
tribunal\footnote{Cas fictif établi pour le cours. Les noms, les dates et les
montants sont fictifs~; les mécanismes techniques et les références
normatives sont exacts.}. Toute l'enquête tient dans une tension~: établir la
vérité \emph{et} pouvoir la prouver.

% ============================================================================
\section{L'Enquêteur n'est pas un Pirate}
\label{sec:p00-posture}
% ============================================================================

\begin{boxdef}
\termecle{Forensique numérique}~: recherche méthodique, documentée et
recevable des traces laissées par un événement sur les supports numériques,
afin d'établir ce qui s'est passé, et de pouvoir le démontrer devant un
tiers. Elle se distingue du piratage par le cadre (autorisation, droit) et la
posture (neutralité, traçabilité).
\end{boxdef}

\begin{tableaucours}{L{4.0cm}L{4.2cm}L{4.0cm}}{%
    \tableauHeader{} &
    \tableauHeader{Le pirate} &
    \tableauHeader{L'enquêteur forensique}
}
But & Accéder, nuire, profiter & Établir et prouver \\
Cadre & Illégal & Autorisé, encadré \\
Trace & Les efface & Les préserve \\
Posture & Partisane & Neutre, impartiale \\
\end{tableaucours}
\vspace{-0.8em}
\begin{center}{\small\itshape Tableau~0.1. Deux usages des mêmes compétences}\end{center}

% ============================================================================
\section{Le Fil Rouge : LiZbiZ et l'Opération TSANAGA}
\label{sec:p00-filrouge}
% ============================================================================

Le cours suit une même enquête, l'\sigle{opération TSANAGA}, menée à Garoua
par le cabinet NYONG SÉCURITÉ pour le compte du parquet, sur les supports
saisis chez M.~NDONGO.

\begin{tableaucours}{L{3.4cm}L{8.8cm}}{%
    \tableauHeader{Acteur} &
    \tableauHeader{Rôle}
}
M.~BIKORO & Expert forensique, cabinet NYONG SÉCURITÉ \\
Mme~NKENG & Analyste forensique \\
Me~TSALA & Avocat~; veille à la recevabilité \\
Le parquet de Garoua & Autorité requérante \\
M.~NDONGO & Mis en cause (ancien employé de LiZbiZ) \\
\end{tableaucours}
\vspace{-0.8em}
\begin{center}{\small\itshape Tableau~0.2. Les acteurs de l'opération TSANAGA}\end{center}

% ============================================================================
\section{Deux Fils : CRO et Chaîne de Possession}
\label{sec:p00-cro}
% ============================================================================

\begin{boxCRO}
\textbf{Le Trilemme CRO.} \termecle{Confidentialité}~: ne saisir et n'exposer
que le nécessaire. \termecle{Fiabilité}~: la trace est-elle fidèle, intègre,
non altérée~? \termecle{Opposabilité}~: pourra-t-on s'en prévaloir devant le
juge~? Chaque acte d'enquête s'arbitre sur ces trois dimensions, et se
\emph{documente}.
\end{boxCRO}

\begin{boxprocedure}
\textbf{La chaîne de possession, dès la saisie}
\begin{enumerate}
    \item Identifier chaque support (numéro de scellé, description).
    \item Consigner qui l'a, depuis quand, pour quoi.
    \item Ne jamais rompre la continuité de la détention.
    \item Calculer les empreintes avant toute analyse.
    \item Joindre la chaîne au dossier.
\end{enumerate}
\end{boxprocedure}

% ============================================================================
\section{Les Questions du Cours}
\label{sec:p00-questions}
% ============================================================================

\begin{boxinfo}
\textbf{Ce que le cours apprend à répondre.}
\begin{itemize}
    \item Qu'est-ce qu'une enquête numérique, et qu'est-ce qui la distingue
          d'un piratage (Partie~I)~?
    \item Qu'est-ce qui rend une preuve recevable (Partie~II)~?
    \item Comment acquérir sans altérer (Partie~III)~?
    \item Comment analyser disque, mémoire, réseau, mobile, nuage
          (Partie~IV)~?
    \item Comment déjouer l'anti-forensique et reconstituer (Partie~V)~?
    \item Comment restituer, par un rapport et devant un tribunal
          (Partie~VI)~?
\end{itemize}
\end{boxinfo}

\separateur

\begin{boxresume}
\textbf{Ce qu'il faut retenir de ce prologue~:}
\begin{itemize}
    \item \textbf{Tout contact laisse une trace}~; l'enquête la cherche et la
          préserve.
    \item \textbf{L'enquêteur n'est pas un pirate}~: cadre et posture le
          distinguent.
    \item \textbf{Le fil rouge}~: LiZbiZ et l'opération TSANAGA à Garoua.
    \item \textbf{CRO et chaîne de possession} encadrent chaque acte.
    \item \textbf{Établir et prouver}~: la double exigence de l'enquête.
\end{itemize}
\end{boxresume}

\begin{boxinfo}
\textbf{Transition.} Commençons par définir la discipline et sa posture~:
c'est l'objet du Chapitre~\ref{chap:forensique}.
\end{boxinfo}
